% BHU PYQ -- Economics: Introductory Microeconomics (2024-25 NEP)
\documentclass[11pt,a4paper]{article}
\usepackage[a4paper,top=2.5cm,bottom=2.5cm,left=2.8cm,right=2.8cm,headheight=14pt]{geometry}
\usepackage{amsmath,amssymb}
\usepackage{fontspec}
\setmainfont{Kohinoor Devanagari}
\usepackage{microtype,fancyhdr,enumitem,hyperref,lastpage,parskip}

\hypersetup{colorlinks=true,urlcolor=black,linkcolor=black,pdftitle={ECOMJ11: Introductory Microeconomics -- BHU NEP 2024-25}}

\pagestyle{fancy}\fancyhf{}
\renewcommand{\headrulewidth}{0.4pt}\renewcommand{\footrulewidth}{0.4pt}
\fancyhead[L]{\small \textit{Banaras Hindu University}}
\fancyhead[R]{\small \textit{ECOMJ11: Microeconomics (2024-25)}}
\fancyfoot[C]{\small Page \thepage\ of \pageref{LastPage}}

\newlist{parts}{enumerate}{2}
\setlist[parts,1]{label=(\alph*),leftmargin=*,itemsep=4pt,topsep=3pt}
\setlist[parts,2]{label=(\roman*),leftmargin=*,itemsep=2pt,topsep=2pt}

\newcommand{\pts}[1]{\hfill \textit{[#1]}}

\begin{document}

\begin{center}
  {\large\bfseries BANARAS HINDU UNIVERSITY}\\[2pt]
  {\normalsize B. A. (Hons.) Semester I Examination, 2024-25 (NEP)}\\[6pt]
  \rule{0.8\linewidth}{0.5pt}\\[6pt]
  {\large\bfseries ECONOMICS}\\[4pt]
  {\large\bfseries Paper Code: ECOMJ11 : Introductory Microeconomics (Major)}\\[4pt]
  {\normalsize Time: 2:30 Hours \hfill Full Marks: 70}\\[2pt]
  {\small REF NO. CE/Conf./D/105}
\end{center}
\rule{\linewidth}{0.4pt}
\smallskip

\noindent \textit{Instructions: This question paper comprises three Sections A, B and C. All questions are compulsory. Write your Roll No.\ at the top immediately on receipt of this question paper.}

\smallskip
\noindent \textit{इस प्रश्न-पत्र में तीन खण्ड अ, ब तथा स हैं। सभी प्रश्न अनिवार्य हैं।}

\bigskip

\section*{SECTION -- A / खण्ड -- अ}
\noindent \textit{Note: Answer each question within 50 words. Each question carries 2 marks.}\pts{5 $\times$ 2 = 10}

\noindent \textit{प्रत्येक प्रश्न का उत्तर 50 शब्दों के भीतर दीजिए। प्रत्येक प्रश्न 2 अंकों का है।}

\begin{enumerate}[label=\textbf{\arabic*.},leftmargin=*,itemsep=8pt]
  \item Explain the following concepts:\pts{5 $\times$ 2 = 10}
  \begin{parts}
    \item Production Possibility Curve / उत्पादन सम्भावना वक्र
    \item Ridge Lines / रिज रेखाएँ
    \item Price Discrimination / कीमत विभेद
    \item Product Differentiation / वस्तु भेद
    \item Break-Even Point / सम-विच्छेद बिन्दु
  \end{parts}
\end{enumerate}

\bigskip

\section*{SECTION -- B / खण्ड -- ब}
\noindent \textit{Note: Write your answer in about 250 words. Each question carries 10 marks.}\pts{3 $\times$ 10 = 30}

\noindent \textit{अपना उत्तर लगभग 250 शब्दों में लिखिए। प्रत्येक प्रश्न 10 अंकों का है।}

\begin{enumerate}[label=\textbf{\arabic*.},leftmargin=*,start=2,itemsep=12pt]

  \item With the help of appropriate diagrams and schedules, discuss the Law of Demand. Also explain why a demand curve slopes downward.\pts{10}
  
  उपयुक्त चित्रों तथा तालिका की सहायता से माँग के नियम की चर्चा कीजिए। यह भी समझाइए कि माँग वक्र नीचे की ओर क्यों झुकता है?

  \begin{center}\textbf{OR / अथवा}\end{center}

  Using compensating variation in income, decompose the price effect into income effect and substitution effect for a normal good.\pts{10}
  
  आय की क्षतिपूरक परिवर्तन का उपयोग करते हुए सामान्य वस्तुओं के लिए कीमत प्रभाव को आय प्रभाव तथा प्रतिस्थापन प्रभाव में विभाजित कीजिए।

  \item Discuss why an average cost curve is U-shaped. Using suitable diagrams, explain the relationship between average cost and marginal cost.\pts{10}
  
  व्याख्या कीजिए कि औसत लागत वक्र U-आकार का क्यों होता है? उपयुक्त चित्रों का उपयोग करते हुए औसत लागत तथा सीमान्त लागत के मध्य सम्बन्ध को समझाइए।

  \begin{center}\textbf{OR / अथवा}\end{center}

  Explain the short-run equilibrium of a firm under perfect competition.\pts{10}
  
  पूर्ण प्रतियोगिता के अंतर्गत एक फर्म के अल्पकालीन साम्य की व्याख्या कीजिए।

  \item Discuss the concept of Consumer's Surplus and explain how it is measured using indifference curve analysis.\pts{10}
  
  उपभोक्ता की बचत की अवधारणा की विवेचना कीजिए तथा उदासीनता वक्र विश्लेषण द्वारा इसके मापन की व्याख्या कीजिए।

  \begin{center}\textbf{OR / अथवा}\end{center}

  Explain the Returns to Scale and distinguish between Increasing, Constant, and Diminishing Returns to Scale.\pts{10}
  
  पैमाने के प्रतिफल की व्याख्या कीजिए तथा बढ़ते, स्थिर एवं घटते पैमाने के प्रतिफल के मध्य अन्तर स्पष्ट कीजिए।

\end{enumerate}

\bigskip

\section*{SECTION -- C / खण्ड -- स}
\noindent \textit{Note: Write your answer in about 500 words. Each question carries 15 marks.}\pts{2 $\times$ 15 = 30}

\noindent \textit{अपना उत्तर लगभग 500 शब्दों में लिखिए। प्रत्येक प्रश्न 15 अंकों का है।}

\begin{enumerate}[label=\textbf{\arabic*.},leftmargin=*,start=5,itemsep=14pt]

  \item Using tables and diagrams, explain the relationship between Total Product and Marginal Product under the Law of Variable Proportions. Also explain the reasons for the occurrence of different stages.\pts{15}
  
  परिवर्तनशील अनुपातों के नियम के अंतर्गत तालिकाओं तथा चित्रों की सहायता से कुल उत्पादन तथा सीमान्त उत्पादन के मध्य सम्बन्ध को समझाइए। विभिन्न अवस्थाओं के उत्पन्न होने के कारणों की भी व्याख्या कीजिए।

  \begin{center}\textbf{OR / अथवा}\end{center}

  What is Producer's Equilibrium? With the help of appropriate diagrams, explain how Producer's Equilibrium is achieved when:
  (i) costs are given, and (ii) output is given.\pts{15}
  
  उत्पादक का साम्य क्या है? उपयुक्त चित्रों की सहायता से समझाइए कि उत्पादक का साम्य किस प्रकार प्राप्त किया जा सकता है जब:
  (i) लागतें दी गई हों, तथा (ii) उत्पादन दिया गया हो।

  \item Explain how price and output are determined under Monopoly. How does a monopolist discriminate prices in different markets?\pts{15}
  
  एकाधिकार के अंतर्गत कीमत तथा उत्पादन का निर्धारण किस प्रकार होता है? एक एकाधिकारी विभिन्न बाज़ारों में कीमत विभेद किस प्रकार करता है?

  \begin{center}\textbf{OR / अथवा}\end{center}

  Discuss the Marginal Productivity Theory of Distribution. Critically evaluate its assumptions and limitations.\pts{15}
  
  वितरण के सीमान्त उत्पादकता सिद्धान्त की विवेचना कीजिए। इसकी मान्यताओं तथा सीमाओं का समालोचनात्मक मूल्यांकन कीजिए।

\end{enumerate}

\end{document}
